\section{Related Work}
\label{sec:related}



% 1. Systems and the law
% 2. Audits of deployed services
% 3. Unlearning and verification
% 4. Memorization and agent memory
% 5. Users and deletion controls
% 6. Users of LLM agents


% \begin{outlinenote}
% Plan: a wide and deep review of papers, published venues only , each verified against the publisher or proceedings page before it is cited.
% \begin{enumerate}[leftmargin=*,nosep]
% \item Data deletion in computer systems: What deletion means formally, and how software systems, databases, storage, websites, and web applications implement data deletion, including where the right to be forgotten has been applied outside language models.
% \item Audits of deployed services, about 10 papers. Technical analyses of whether websites, mobile apps, and other platforms delete data when asked.
% \item Deletion in AI: Machine unlearning methods, attacks on unlearning, how unlearning is verified, memorization of training data, and whether unlearning can be made usable by users.
% \item Users:. How users perceive privacy and memory in AI systems, how usable and effective deletion and privacy controls are, and how users verify that their data was removed.
% \item Gaps. Prior work audits services, studies unlearning inside models, and studies how users use controls. No work tests what happens after a deletion request to deployed LLM agents, or compares the outcome with what users expect and how they try to verify it.
% \end{enumerate}
% \end{outlinenote}


\para{Data deletion effectiveness and usability}
%
Prior research has examined the effectiveness and usability of data deletion across systems (mostly in domains other than LLM agents).
%
For example, secure deletion research shows that removing a file or issuing a deletion command does not necessarily eliminate the underlying data~\cite{reardon2013sok}.
%
Similar challenges arise in databases, where physical removal may be delayed~\cite{sarkar2020lethe,shastri2020understanding}, and in backups that fail to support the right to be forgotten~\cite{politou2018backups}.
%
Studies of deployed services have also identified gaps between deletion controls and actual outcomes.
%
Habib et al.~\cite{habib2019empirical} identified limitations in how websites present and explain deletion and opt-out choices, while other studies found that online services and mobile applications may retain personal information after deletion requests or account deletion~\cite{rupp2022leave,santhanam2022scraping,yan2025sign}.
%
These findings highlight the need to evaluate deletion outcomes rather than assume that available controls ensure data removal.

Beyond deletion effectiveness, usable privacy research further shows that users struggle to locate and exercise deletion controls~\cite{ramokapane2017stupid,habib2020scavenger,bumiller2026zombies}, and that their expectations may differ from what services provide~\cite{rao2016expecting}.
%
Minaei et al.~\cite{minaei2022deletion} also identified limitations in deletion privacy through users' experiences and expectations.
%
These usability studies are loosely related to our work, as they may affect how users choose, understand, and verify data deletion.

\para{Unlearning in machine learning}
%
Machine unlearning aims to remove the influence of selected training data from learned models.
%
Prior work has developed formal definitions~\cite{garg2020formalizing}, methods that avoid complete retraining~\cite{cao2015towards,bourtoule2021machine}, and certified removal guarantees under specific assumptions~\cite{guo2020certified}, with surveys summarizing these advances~\cite{xu2023machine}.
%
However, verification remains challenging~\cite{thudi2022necessity,xue2026towards}.
%
Evaluations of language-model unlearning reveal persistent privacy risks, including information leakage after unlearning~\cite{shi2025muse,patil2024sensitive}.
%
In-context unlearning can also modify model behavior without removing information from model parameters~\cite{pawelczyk2024incontext}, illustrating why non-retrieval does not necessarily establish erasure.
%
Related research shows that language models can memorize and reproduce training data~\cite{carlini2019secret,carlini2021extracting,carlini2023quantifying,lukas2023analyzing}.
%
Controlled insertion and extraction of synthetic information provides a useful way to investigate memorization~\cite{carlini2019secret}.
%
Building on this principle, our study examines a different setting: whether user-provided information remains retrievable after deletion through the controls of deployed LLM agents, rather than unlearning from model parameters.

\para{Data deletion in LLM agents}
%
LLM agents introduce additional privacy challenges by retaining user interactions and derived information.
%
As shown in past research, generative agents can store past interactions and generate reflections~\cite{park2023generative}, while agent memory can expose private information to extraction attacks~\cite{wang2025unveiling}, raising questions about whether deletion controls cover retained and derived data.
%
Recent studies have examined privacy controls and user perceptions in conversational LLM platforms.
%
For example, Li et al.~\cite{li2026privacy} studied six platforms, identifying challenges in accessing, editing, deleting, and sharing conversational data, including ambiguity in natural-language privacy controls.
%
Other work examined users' privacy concerns, disclosure decisions, and expectations~\cite{zhang2024fair,malki2025hoovered,kwesi2026impact}, while studies of ChatGPT memory and retrieval-augmented memory highlighted tensions between personalization and privacy~\cite{chen2026relational,zhang2025understanding}.
%
Despite these advances, limited empirical evidence connects users' deletion expectations with the observed post-deletion behavior of commercial LLM agents.
%
Our study addresses this gap by combining a user survey with controlled post-deletion recall experiments to evaluate information retrievability after different deletion methods and compare observed outcomes with user expectations and platform documentation.